\documentclass[../DoAn.tex]{subfiles}
\begin{document}

\section{Triển khai mô hình Chatbot TinyLLaMA}

Phần này trình bày chi tiết toàn bộ quá trình triển khai hệ thống chatbot sử dụng mô hình ngôn ngữ nhỏ gọn TinyLLaMA. Mục tiêu chính là xây dựng một chatbot nhẹ, hiệu quả, có khả năng sinh phản hồi cho các câu hỏi thực tế từ sinh viên dựa trên dữ liệu email đã được xử lý. Hệ thống được phát triển trên môi trường Python với sự hỗ trợ của các thư viện hiện đại như \texttt{transformers}, \texttt{peft}, \texttt{datasets}, \texttt{accelerate}, và \texttt{bitsandbytes}.

Các bước triển khai bao gồm:
\begin{itemize}
    \item Chuẩn bị dữ liệu huấn luyện từ email.
    \item Định dạng prompt theo cấu trúc hội thoại.
    \item Tải mô hình và tokenizer từ HuggingFace.
    \item Áp dụng kỹ thuật LoRA để fine-tune mô hình.
    \item Token hóa dữ liệu và tạo tập huấn luyện.
    \item Huấn luyện mô hình bằng \texttt{Trainer}.
    \item Sinh phản hồi từ mô hình đã fine-tune.
\end{itemize}

\section{Chuẩn bị dữ liệu huấn luyện}

Tập dữ liệu huấn luyện được lưu trong tệp \texttt{qa\_dataset\_cleaned.jsonl}, bao gồm các cặp hỏi – đáp được trích xuất từ nội dung email giữa sinh viên và giảng viên. Mỗi dòng trong tệp là một đối tượng JSON có định dạng như sau:

\begin{verbatim}
{
  "instruction": "Câu hỏi của sinh viên",
  "output": "Câu trả lời của giảng viên"
}
\end{verbatim}

Để phù hợp với mô hình sinh ngôn ngữ, mỗi cặp được định dạng lại theo cấu trúc hội thoại dạng prompt:

\begin{lstlisting}
### User:
[instruction]

### Assistant:
[output]
\end{lstlisting}

Việc định dạng này giúp mô hình hiểu rõ vai trò của từng bên trong hội thoại, từ đó sinh ra phản hồi tự nhiên và chính xác hơn.

\textbf{Nạp dữ liệu:}
\begin{verbatim}
with open("/content/qa_dataset_cleaned.jsonl", "r", 
          encoding="utf-8") as f:
    data = [json.loads(line) for line in f]
\end{verbatim}

\textbf{Chuyển đổi thành Dataset:}
\begin{verbatim}
dataset = Dataset.from_list(data)
\end{verbatim}

\textbf{Định dạng prompt:}
\begin{verbatim}
def format_example(example):
    return {
        "text": f"### User:\n{example['instruction']}"
                f"\n\n### Assistant:\n{example['output']}"
    }
\end{verbatim}

\begin{verbatim}
dataset = dataset.map(format_example, 
                     remove_columns=["instruction", "output"])
\end{verbatim}

Mô hình được sử dụng là \texttt{TinyLLaMA-1.1B-Chat-v1.0} với 1.1 tỷ tham số – một mô hình nhẹ, thích hợp triển khai trên máy tính cá nhân. Tokenizer được cấu hình để sử dụng phiên bản nhanh và đặt token \texttt{pad} trùng với token kết thúc chuỗi để đồng nhất xử lý chuỗi.

\begin{lstlisting}
from transformers import AutoTokenizer, AutoModelForCausalLM

model_id = "TinyLlama/TinyLlama-1.1B-Chat-v1.0"
tokenizer = AutoTokenizer.from_pretrained(model_id, use_fast=True)
tokenizer.pad_token = tokenizer.eos_token

model = AutoModelForCausalLM.from_pretrained(
    model_id,
    load_in_8bit=True,
    device_map="auto"
)
\end{lstlisting}

\section{Áp dụng kỹ thuật LoRA với PEFT}

Để fine-tune mô hình hiệu quả, đề tài sử dụng kỹ thuật LoRA (Low-Rank Adaptation) qua thư viện \texttt{peft}. LoRA giúp giảm số lượng tham số cần huấn luyện, phù hợp với các mô hình lớn hoặc phần cứng giới hạn.

\begin{itemize}
    \item \textbf{Rank (r):} 16
    \item \textbf{Hệ số khuếch đại (lora\_alpha):} 32
    \item \textbf{Target modules:} \texttt{["q\_proj", "v\_proj"]}
    \item \textbf{Dropout:} 0.05
    \item \textbf{Bias:} none
    \item \textbf{Task type:} \texttt{CAUSAL\_LM}
\end{itemize}

\begin{lstlisting}
from peft import get_peft_model, LoraConfig, TaskType

lora_config = LoraConfig(
    r=16,
    lora_alpha=32,
    target_modules=["q_proj", "v_proj"],
    lora_dropout=0.05,
    bias="none",
    task_type=TaskType.CAUSAL_LM
)

model = get_peft_model(model, lora_config)
model.print_trainable_parameters()
\end{lstlisting}

\section{Token hóa và chuẩn bị tập huấn luyện}

Sau khi định dạng prompt, dữ liệu được token hóa để chuẩn bị huấn luyện. Tokenizer chuyển văn bản thành chuỗi số với các thiết lập sau:

\begin{itemize}
    \item \texttt{truncation=True}: Cắt bớt nếu vượt quá 512 token.
    \item \texttt{padding="max\_length"}: Đệm cho đủ độ dài cố định.
    \item \texttt{max\_length=512}: Giới hạn độ dài tối đa.
\end{itemize}

\begin{lstlisting}
def tokenize(example):
    return tokenizer(
        example["text"],
        truncation=True,
        padding="max_length",
        max_length=512
    )

tokenized_dataset = dataset.map(tokenize, batched=True)
\end{lstlisting}

\section{Cấu hình và thực thi huấn luyện}

Mô hình được huấn luyện bằng \texttt{Trainer} trong thư viện \texttt{transformers}. Cấu hình chi tiết như sau:

\begin{itemize}
    \item \texttt{batch\_size}: 4 mẫu mỗi thiết bị
    \item \texttt{gradient\_accumulation\_steps}: 2
    \item \texttt{epoch}: 3
    \item \texttt{learning\_rate}: 2e-4
    \item \texttt{fp16}: bật huấn luyện nửa chính xác
    \item \texttt{logging/save steps}: 10/200
    \item \texttt{output\_dir}: \texttt{./tinyllama-lora}
\end{itemize}

\begin{lstlisting}
from transformers import TrainingArguments, Trainer, DataCollatorForLanguageModeling

data_collator = DataCollatorForLanguageModeling(
    tokenizer=tokenizer,
    mlm=False
)

training_args = TrainingArguments(
    output_dir="./tinyllama-lora",
    per_device_train_batch_size=4,
    gradient_accumulation_steps=2,
    num_train_epochs=3,
    learning_rate=2e-4,
    fp16=True,
    logging_steps=10,
    save_steps=200,
    save_total_limit=1,
    report_to="none"
)

trainer = Trainer(
    model=model,
    args=training_args,
    train_dataset=tokenized_dataset,
    tokenizer=tokenizer,
    data_collator=data_collator
)

trainer.train()

model.save_pretrained("./tinyllama-lora")
tokenizer.save_pretrained("./tinyllama-lora")
\end{lstlisting}

\section{Sinh phản hồi từ mô hình đã huấn luyện}

Sau huấn luyện, mô hình được nạp lại để sinh phản hồi từ các câu hỏi thực tế. Pipeline được cấu hình để sinh văn bản mạch lạc và đa dạng:

\begin{itemize}
    \item \texttt{max\_new\_tokens=1000}
    \item \texttt{do\_sample=True}
    \item \texttt{temperature=0.7}
\end{itemize}

\begin{lstlisting}
from transformers import pipeline, AutoTokenizer, AutoModelForCausalLM
from peft import PeftModel

tokenizer = AutoTokenizer.from_pretrained("./tinyllama-lora")

base_model = AutoModelForCausalLM.from_pretrained(
    "TinyLlama/TinyLlama-1.1B-Chat-v1.0",
    load_in_8bit=True,
    device_map="auto"
)

model = PeftModel.from_pretrained(base_model, "./tinyllama-lora")
model.eval()

pipe = pipeline("text-generation", model=model, tokenizer=tokenizer)

prompt = "### User:\nCho em hỏi ngành CNTT Việt Nhật cần đầu ra tiếng Anh không ạ\n\n### Assistant:\n"

result = pipe(prompt, max_new_tokens=1000, do_sample=True, temperature=0.7)
print(result[0]["generated_text"])
\end{lstlisting}

\section{Kết luận chương}

Chương này đã trình bày chi tiết toàn bộ quy trình triển khai chatbot sử dụng mô hình TinyLLaMA kết hợp kỹ thuật fine-tune LoRA. Hệ thống đạt được hiệu quả cao với chi phí tài nguyên thấp, phù hợp cho triển khai thực tế trên thiết bị cá nhân. Trong chương tiếp theo, mô hình sẽ được đánh giá về chất lượng phản hồi, độ chính xác và tính khả thi trong ứng dụng thực tế.

\end{document}
